% Lösungen Kurvenuntersuchung (zusammenbau v0.4, Heft)
\einheitenkopf{Kurvenuntersuchung}

\erg{16}{a) steigt – die Tangente in $P$ geht nach rechts oben. \quad b) $f$ steigt auf diesem Intervall – für $x \ge 4$ ist $f'(x) = 2x - 6 \ge 2$, also positiv. \quad c) $f'(x) = -3x^2 + 6x + 9 = 0$ für $x_1 = -1$ und $x_2 = 3$; $f'(-2) = -15 < 0$: $f$ fällt auf $]-\infty; -1]$; $f'(0) = 9 > 0$: $f$ steigt auf $[-1; 3]$; $f'(4) = -15 < 0$: $f$ fällt auf $[3; \infty[$.}
\erg{17}{a) $f'(x) = x^2 \ge 0$ für alle $x$; $f'$ ist nur an der einzelnen Stelle $x = 0$ null, also ist der Graph monoton steigend. \quad b) $f'(x) = 2\,\mathrm{e}^{2x} > 0$ für alle $x$, weil $\mathrm{e}^{2x}$ immer positiv ist; also ist der Graph streng monoton steigend. \quad c) $d(x) = \frac{1}{4}x^3 - 3x + 5$; $d'(x) = \frac{3}{4}x^2 - 3 = 0$ für $x = 2$ (im Intervall); $d$ fällt auf $[0; 2]$ und steigt auf $[2; 4]$; $d(2) = 1$, $d(0) = 5$, $d(4) = 9$; Wertebereich $[1; 9]$ – das Maximum liegt am rechten Rand. \quad d) $d(x) = \frac{1}{2}x^3 - 6x + 10$; $d'(x) = \frac{3}{2}x^2 - 6 = 0$ für $x = 2$ (im Intervall); $d$ fällt auf $[0; 2]$ und steigt auf $[2; 3]$; $d(2) = 2$, $d(0) = 10$, $d(3) = 5{,}5$; Wertebereich $[2; 10]$ – das Maximum liegt am linken Rand.}
\erg{18}{a) Stelle gegeben – nachweisen: $x = 2$ einsetzen. \quad b) $f'(x) = 3x^2 - 3x - 6 = 0 \Leftrightarrow x^2 - x - 2 = 0$, also $x_1 = 2$; $x_2 = -1$ \quad c) $f'(x) = 4x^3 - 16x = 4x(x^2 - 4) = 0$ für $x_1 = 0$, $x_2 = -2$, $x_3 = 2$; $f''(x) = 12x^2 - 16$; $f''(0) = -16 < 0$: Hochpunkt $H(0 | 0)$; $f''(\pm 2) = 32 > 0$: Tiefpunkte $T_1(-2 | -16)$ und $T_2(2 | -16)$}
\erg{19}{a) Den Scheitel der Parabel bestimmen und von seiner $y$-Koordinate die Lkw-Höhe abziehen: $p'(x) = -0{,}5x + 2 = 0$ für $x = 4$; $p(4) = 5$; Abstand $5 - 4 = 1$ m. \quad b) $f'(x) = -\mathrm{e}^{x} + (3 - x)\,\mathrm{e}^{x} = (2 - x)\,\mathrm{e}^{x} = 0$; da $\mathrm{e}^{x} > 0$, ist $x = 2$; $f''(x) = (1 - x)\,\mathrm{e}^{x}$, $f''(2) < 0$; $H(2 | \mathrm{e}^{2})$, also $H(2 | 7{,}39)$ \quad c) $c'(t) = 6\,\mathrm{e}^{-0{,}2t} - 1{,}2t\,\mathrm{e}^{-0{,}2t} = (6 - 1{,}2t)\,\mathrm{e}^{-0{,}2t} = 0$; da der e-Faktor nie null ist, ist $t = 5$; $c(5) = 30\,\mathrm{e}^{-1} \approx 11{,}04$ mg pro Liter nach $5$ Stunden \quad d) $f'(x) = 2(x - 2)\,\mathrm{e}^{-0{,}5x} - 0{,}5(x - 2)^2\,\mathrm{e}^{-0{,}5x} = (-0{,}5x^2 + 4x - 6)\,\mathrm{e}^{-0{,}5x} = 0 \Leftrightarrow x^2 - 8x + 12 = 0$, also $x_1 = 2$, $x_2 = 6$; $f''(2) = 2\,\mathrm{e}^{-1} > 0$: $T(2 | 0)$; $f''(6) = -2\,\mathrm{e}^{-3} < 0$: $H(6 | 16\,\mathrm{e}^{-3})$, also $H(6 | 0{,}80)$ \quad e) $f'(x) = (-20x + 30)\,\mathrm{e}^{-x} - (-10x^2 + 30x - 10)\,\mathrm{e}^{-x} = (10x^2 - 50x + 40)\,\mathrm{e}^{-x} = 0 \Leftrightarrow x^2 - 5x + 4 = 0$, also $x_1 = 1$, $x_2 = 4$; $f''(1) = -30\,\mathrm{e}^{-1} < 0$: $H(1 | 10\,\mathrm{e}^{-1})$, also $H(1 | 3{,}68)$; $f''(4) = 30\,\mathrm{e}^{-4} > 0$: $T(4 | -50\,\mathrm{e}^{-4})$, also $T(4 | -0{,}92)$}
\erg{20}{a) Stelle gegeben – nachweisen: $x = -1$ in $f'$ und $f''$ einsetzen. \quad b) $f'(x) = 6x^2 - 12x$; $f'(2) = 24 - 24 = 0$ \quad c) $f'(x) = 4x^3 - 32$, $f'(2) = 0$; $f'(0) = -32 < 0$, $f'(3) = 76 > 0$, und $f'$ hat nur die Nullstelle $2$: Wechsel von $-$ nach $+$, also Tiefpunkt $T(2 | -48)$.}
\erg{21}{a) $f'(x) = 3x^2 - 27$, $f'(-3) = 27 - 27 = 0$; $f''(x) = 6x$, $f''(-3) = -18 < 0$, also Hochpunkt. \quad b) $f'(x) = 3x^2 - 6x = 3x(x - 2)$, $f'(0) = f'(2) = 0$; $f'(-1) = 9 > 0$, $f'(1) = -3 < 0$, $f'(3) = 9 > 0$: bei $0$ Wechsel von $+$ nach $-$ (Hochstelle), bei $2$ von $-$ nach $+$ (Tiefstelle). \quad c) $N(-3 | 0)$ und $S_y(0 | 3)$; $f'(x) = 0$ nur für $x = -2$, da $\mathrm{e}^{-x} > 0$; für $x < -2$ ist $f'(x) > 0$, für $x > -2$ ist $f'(x) < 0$: Wechsel von $+$ nach $-$, also Hochpunkt $H(-2 | \mathrm{e}^{2})$, also $H(-2 | 7{,}39)$.}
\erg{22}{a) $p'(x) = -2x - 4 = 0$ für $x = -2$, Scheitel $S(-2 | 5)$; Abstand $d = \sqrt{(-2 + 1)^2 + (5 - 4)^2} = \sqrt{2} \approx 1{,}41 > 1$.}
\erg{23}{a) $p'(x) = -2x + 6 = 0$ für $x = 3$, Scheitel $S(3 | 5{,}5)$; Abstand $d = \sqrt{(3 - 2)^2 + (5{,}5 - 6)^2} = \sqrt{1{,}25} \approx 1{,}12 > 0{,}5$. Die Differenz der $y$-Werte allein wäre nur $0{,}5$.}
\erg{24}{a) $f(-2) = -4 + 12 = 8$, also $H(-2 | 8)$; wegen der Punktsymmetrie $T(2 | -8)$; Abstand $d = \sqrt{4^2 + 16^2} = \sqrt{272} \approx 16{,}49$}
\erg{25}{a) $f(-1) = -2 + 6 = 4$, also $H(-1 | 4)$; wegen der Punktsymmetrie $T(1 | -4)$; Abstand $d = \sqrt{2^2 + 8^2} = \sqrt{68} \approx 8{,}25$}
\erg{26}{a) $f'(x) = -\frac{3}{8}x^2 + 1{,}5 = 0 \Leftrightarrow x^2 = 4$; $f''(x) = -\frac{3}{4}x$: $T(-2 | -2)$, $H(2 | 2)$; Steigung $\frac{2 - (-2)}{2 - (-2)} = 1$ und $H$ liegt auf $y = x$: die Gerade ist $y = x$.}
\erg{27}{a) $f'(x) = \frac{3}{8}x^2 - 1{,}5 = 0 \Leftrightarrow x^2 = 4$; $f''(x) = \frac{3}{4}x$: $H(-2 | 2)$, $T(2 | -2)$; Steigung $\frac{-2 - 2}{2 - (-2)} = -1$ und $H$ liegt auf $y = -x$: die Gerade ist $y = -x$.}
\erg{28}{a) $x^2 + 4x - 5 = 0$: $x_1 = -5$, $x_2 = 1$; Extremstelle in der Mitte: $x = -2$}
\erg{29}{a) $x^2 - 2x - 3 = 0$: $x_1 = -1$, $x_2 = 3$; Extremstelle in der Mitte der Nullstellen: $x = 1$}
\erg{30}{a) $g_o(x) - f(x) = \frac{1}{16}x(x - 6)^2 \ge 0$ für $x \ge 0$; $d(x) = f(x) - g_u(x) = -\frac{1}{16}x^3 + 0{,}75x^2 - 2{,}25x + 4$, $d'(x) = -\frac{3}{16}(x - 2)(x - 6) = 0$ für $x = 2$ (Minimum, $d(2) = 2$) und $x = 6$ (Maximum, $d(6) = 4$); Randwerte $d(0) = 4$, $d(8) = 2$; also $d > 0$. Kleinster Abstand im Inneren bei $x = 2$, er beträgt $2$.}
\erg{31}{a) $g_o(x) - f(x) = 0{,}25x(x - 3)^2 \ge 0$ für $x \ge 0$; $d(x) = f(x) - g_u(x) = -0{,}25x^3 + 1{,}5x^2 - 2{,}25x + 3$, $d'(x) = -0{,}75(x - 1)(x - 3) = 0$ für $x = 1$ (Minimum, $d(1) = 2$) und $x = 3$ (Maximum, $d(3) = 3$); Randwerte $d(0) = 3$, $d(4) = 2$; also $d > 0$ auf $[0; 4]$. Kleinster Abstand im Inneren bei $x = 1$, er beträgt $2$.}
\erg{32}{a) $f'''(x_0) \ne 0$ – der Wert $f'''(1) = 6$ ist gegeben. \quad b) $f'(x) = 3x^2 - 12x + 4$; $f''(x) = 6x - 12$; $f'''(x) = 6$ \quad c) $f'(x) = (2x - x^2)\,\mathrm{e}^{-x}$; $f''(x) = (2 - 2x)\,\mathrm{e}^{-x} - (2x - x^2)\,\mathrm{e}^{-x} = (x^2 - 4x + 2)\,\mathrm{e}^{-x} = 0$, also $x = 2 \pm \sqrt{2}$; $W_1(0{,}59 | 0{,}19)$, $W_2(3{,}41 | 0{,}38)$}
\erg{33}{a) $f''(x) = \frac{1}{4}(6x - 18)$, $f''(3) = 0$; $f'''(3) = 1{,}5 \ne 0$; $f(3) = 0$; $f'(x) = \frac{1}{4}(3x^2 - 18x + 18)$, $f'(3) = -2{,}25$; Wendetangente $t(x) = -2{,}25(x - 3) = -2{,}25x + 6{,}75$ \quad b) $f''(x) = (2x + 2)\,\mathrm{e}^{x} + (x^2 + 2x - 1)\,\mathrm{e}^{x} = (x^2 + 4x + 1)\,\mathrm{e}^{x} = 0 \Leftrightarrow x^2 + 4x + 1 = 0$, also $x = -2 \pm \sqrt{3}$; $W_1(-3{,}73 | 0{,}31)$, $W_2(-0{,}27 | -0{,}71)$ \quad c) Tiefpunkt bei $x = 2$: $T(2 | -1{,}47)$ im IV. Quadranten, dort ist der Graph linksgekrümmt; rechts davon steigt er und nähert sich von unten der $x$-Achse, ohne sie zu erreichen – dafür muss er in eine Rechtskrümmung übergehen, dort liegt ein Wendepunkt. $g''$ muss man nicht berechnen. \quad d) Tiefpunkt bei $x = 4$: $T(4 | -5{,}89)$ im IV. Quadranten, dort ist der Graph linksgekrümmt; rechts davon steigt er und nähert sich von unten der $x$-Achse – dafür muss er in eine Rechtskrümmung übergehen, dort liegt ein Wendepunkt.}

33c)

\begin{ksys}[xmin=-1,xmax=10,ymin=-2,ymax=1,klein] \funktionab{-2*\x*exp(-0.5*\x)}{g}{-0.5}{10} \tiefpunkt{2}{-1.47}{T} \end{ksys}


33d)

\begin{ksys}[xmin=-1,xmax=16,ymin=-7,ymax=1,xstep=2,klein] \funktionab{-4*\x*exp(-0.25*\x)}{g}{-0.2}{16} \tiefpunkt{4}{-5.89}{T} \end{ksys}

\erg{34}{a) $f''(x) = 12x^2 - 12 = 0$ für $x = \pm 1$, $f'''(\pm 1) \ne 0$: $W_1(-1 | -5)$, $W_2(1 | -5)$; $g(x) = -5$; die Parallele berührt den Bogen im Hochpunkt $(0 | 0)$: $y = 0$.}

34a)

\begin{ksys}[xmin=-2,xmax=2,ymin=-6,ymax=2,klein] \funktionab{\x^4-6*\x^2}{f}{-1}{1} \gerade{0}{-5}{g} \gerade{0}{0}{p} \end{ksys}

\erg{35}{a) $f''(x) = -6x^2 + 12x = -6x(x - 2)$, Vorzeichenwechsel bei $0$ und $2$: $W_1(0 | 0)$, $W_2(2 | 8)$; $g(x) = 4x$; die Parallele berührt den Bogen dazwischen: $f'(x) = -2x^3 + 6x^2 = 4$ für $x = 1$, Berührpunkt $(1 | 1{,}5)$, Gerade $y = 4x - 2{,}5$.}

35a)

\begin{ksys}[xmin=-1,xmax=3,ymin=-3,ymax=9,ystep=2,klein] \funktionab{-0.5*\x^4+2*\x^3}{f}{0}{2} \gerade{4}{0}{g} \gerade{4}{-2.5}{p} \end{ksys}

\erg{36}{a) $f''(x) = -6x + 6$, $f''(1) = 0$; $f'(x) = -3x^2 + 6x - 4$, $f'(1) = -1$; $\mathrm{tan}\,\alpha = 1$, also $\alpha = 45°$.}
\erg{37}{a) $f''(x) = 6x + 6$, $f''(-1) = 0$; $f'(x) = 3x^2 + 6x + 2$, $f'(-1) = -1$; $\mathrm{tan}\,\alpha = 1$, also $\alpha = 45°$.}
\erg{38}{a) $A$: plus, $B$: null, $C$: minus \quad b) Nullstellen von $f'$: $x = -1$ und $x = 1$ (Hoch- und Tiefpunkt); $f' > 0$ für $x < -1$, $f' < 0$ zwischen $-1$ und $1$, $f' > 0$ für $x > 1$ \quad c) $f'(1) = 1{,}5$; Wendestelle $x = 2$ – dort ist der Graph am steilsten.}
\erg{39}{a) $f$: Tiefpunkt $T(0 | -2)$, Hochpunkt $H(4 | 6)$, Randwerte $f(-1) = -0{,}25$ und $f(5) = 4{,}25$; $f'$: nach unten geöffnete Parabel mit den Nullstellen $0$ und $4$ (Extremstellen von $f$) und dem Scheitel $(2 | 3)$. \quad b) I wahr: die am stärksten fallende Tangente liegt im Wendepunkt, $f'(5) = -1$; die Gerade $AW$ hat die Steigung $m = 1$, und $1 \cdot (-1) = -1$. II falsch: $f'$ hat bei $5$ sein Minimum $-1$, die Gerade $y = -1$ berührt den Graphen von $f'$. \quad c) I wahr: $f'(2) = -1$ ist die kleinste Steigung; die Gerade $AW$ hat die Steigung $m = 1$, und $1 \cdot (-1) = -1$. II wahr: der kleinste Wert von $f'$ ist $-1$, Geraden darunter treffen den Graphen von $f'$ nicht. \quad d) Graph durch $(0 | 1)$, steigend bis $H(1 | 5)$, fallend über $W(2 | 3)$ bis $T(3 | 1)$, steigend bis $(4 | 5)$. Die Aussage ist falsch: die Wendetangente hat nur $W$ mit dem Graphen gemeinsam; die Tangenten in $H$ und $T$ haben dagegen zwei gemeinsame Punkte. \quad e) Graph durch $(-3 | 5)$, fallend bis $T(-2 | 1)$, steigend über $W(-1 | 3)$ bis $H(0 | 5)$, fallend bis $(1 | 1)$. Die Aussage ist falsch: die Tangente $y = 5$ trifft den Graphen auch in $(-3 | 5)$.}

39a)

\begin{ksys}[xmin=-1,xmax=5,ymin=-4,ymax=7,klein] \funktionab{-0.25*\x^3+1.5*\x^2-2}{f}{-1}{5} \funktionab{-0.75*\x^2+3*\x}{f'}{-1}{5} \end{ksys}


39d)

\begin{ksys}[xmin=-1,xmax=5,ymin=-1,ymax=6,klein] \funktionab{\x^3-6*\x^2+9*\x+1}{f}{0}{4} \hochpunkt{1}{5}{H} \tiefpunkt{3}{1}{T} \wendepunkt{2}{3}{W} \end{ksys}


39e)

\begin{ksys}[xmin=-4,xmax=2,ymin=-1,ymax=6,klein] \funktionab{-\x^3-3*\x^2+5}{f}{-3}{1} \hochpunkt{0}{5}{H} \tiefpunkt{-2}{1}{T} \wendepunkt{-1}{3}{W} \end{ksys}

\erg{40}{a) stärkste Änderung: Extremum von $f'$, Wendepunkt – gefragt ist die größte Wachstumsrate. \quad b) Die Temperatur fällt von $90$ °C zunächst schnell, dann immer langsamer und nähert sich $20$ °C, der Raumtemperatur; sie sinkt nicht darunter. \quad c) $25$ cm – die steilste Stelle ist der Wendepunkt bei $t = 10$.}
\erg{41}{a) Bei $t = 10$ steigt der Ballon am schnellsten; danach steigt er weiter, aber langsamer – der Wendepunkt ist kein Richtungswechsel. \quad b) $f'(x) = -0{,}012x^2 + 0{,}18x$, $f''(x) = -0{,}024x + 0{,}18 = 0$ für $x = 7{,}5$; $f(7{,}5) \approx 5{,}38$; $f'(7{,}5) = 0{,}675$; $\mathrm{tan}\,\alpha = 0{,}675$, also $\alpha \approx 34{,}0° > 30°$: die Bedingung ist nicht erfüllt. \quad c) Gerade durch $(0 | 0)$ und $(9 | 216)$; sie schneidet den Graphen von $K$ bei $x = 2$ und $x = 8$ und liegt dazwischen oberhalb: Gewinn für Mengen zwischen $2$ m³ und $8$ m³ – links von $2$ nicht, dort sind die Kosten höher. \quad d) $c''(t) = 0 \Leftrightarrow 0{,}24t - 2{,}4 = 0 \Leftrightarrow t = 10$; $c(10) = 60\,\mathrm{e}^{-2} \approx 8{,}12$ mg pro Liter; $c'(t) = (6 - 1{,}2t)\,\mathrm{e}^{-0{,}2t}$, $c'(10) = -6\,\mathrm{e}^{-2} \approx -0{,}81$ mg pro Liter und Stunde. \quad e) $h''(x) = 0{,}0024x^2 - 0{,}24 = 0 \Leftrightarrow x^2 = 100$, also $x = 10$ ($x = -10$ liegt nicht auf dem Hang); $h(10) = 10$; $K(10 | 10)$ mit dem Gefälle $h'(10) = -1{,}6$. \quad f) $g(-3) = -3 + 4{,}5 = 1{,}5$: Höhe $2 \cdot 1{,}5 = 3$ m; die Halbkreise haben den Radius $1{,}5$ m, Breite $2 \cdot (3 + 1{,}5) = 9$ m. Wegen $3 \le 3{,}2$ und $9 \le 9{,}2$ passt die Bahn. Ohne Halbkreise wäre die Breite nur $6$ m.}

41c)

\begin{ksys}[xmin=0,xmax=9,ymin=0,ymax=300,xstep=1,ystep=50,klein] \funktionab{\x^3-9*\x^2+30*\x+16}{K}{0}{9} \gerade{24}{0}{E} \end{ksys}

\erg{42}{a) $d(t)$ ist der Höhenunterschied, um den $A$ größer ist als $B$. Nach etwa $4{,}2$ Wochen ist dieser Unterschied am größten; er beträgt dann $d(4{,}2) \approx 2{,}6$ cm – eine Länge, keine Zeit. \quad b) $d(t)$ ist der Vorsprung von Zug $A$ in km. Nach etwa $12{,}5$ Minuten ist der Vorsprung am größten; er beträgt dann $d(12{,}5) \approx 1{,}8$ km – ein Weg, kein Zeitvorsprung.}
